\documentclass[11pt]{article}
\usepackage{amsmath,amsthm,mathtools,tasks,tcolorbox,multicol,siunitx,fancyhdr,lastpage}
\pagestyle{fancy}
\fancyhf{}
\lhead{Trường THPT Lê Lợi}
\rhead{Đề số 01}
\cfoot{Trang \thepage/\pageref{LastPage}}
\definecolor{xanh}{RGB}{0,102,204}
\newcommand*{\mycmd}[2][Mặc định]{\textbf{#1}: #2}
\NewDocumentCommand{\diemso}{s m O{điểm}}{\IfBooleanTF{#1}{(#2 #3)}{[#2 #3]}}
\DeclarePairedDelimiter{\ceilx}{\lceil}{\rceil}
\def\vect#1{\overrightarrow{#1}}
\def\khoang#1;#2;{\left[#1;#2\right]}
\let\epsi\varepsilon
\theoremstyle{definition}
\newtheorem{dn}{Định nghĩa}[section]
\newtheorem*{nx}{Nhận xét}
\begin{document}
\tableofcontents
\section{Các macro lạ}
\mycmd{giá trị} và \mycmd[Tùy chọn]{khác}. Điểm: \diemso{2} hoặc \diemso*{3}[đ].
Làm tròn $\ceilx{x}$, vectơ $\vect{MN}$, khoảng $\khoang a;b;$, $\epsi>0$.
\tieudelaxa{Văn bản trong lệnh lạ} rồi \lenhkhongbiet\ tiếp tục \caidat{2cm}{nội dung giữ lại}.

\begin{dn}\label{dn:1}
Hàm số $f$ liên tục tại $x_0$ nếu $\lim\limits_{x\to x_0} f(x)=f(x_0)$.
\end{dn}
\begin{nx}
Theo Định nghĩa~\ref{dn:1}, \textit{mọi đa thức} đều liên tục.
\end{nx}
\begin{proof}
Hiển nhiên.
\end{proof}

\subsection{Toán học nâng cao}
\begin{equation}\label{eq:main}
\sum_{k=1}^{n} \binom{n}{k} = 2^n - 1, \quad \operatorname*{argmax}_{x\in X} g(x)
\end{equation}
\begin{gather*}
A = \begin{pmatrix} 1 & 2 \\ 3 & 4 \end{pmatrix}, \quad \det A = \begin{vmatrix} 1 & 2\\ 3 & 4\end{vmatrix} = -2
\end{gather*}
\[
f(x)=\begin{cases} x^2 & \text{nếu } x\ge 0\\ -x & \text{nếu } x<0 \end{cases}, \qquad
\underbrace{1+1+\cdots+1}_{n \text{ lần}} = n, \qquad \mathcal{F}, \mathfrak{g}, \mathbb{Z}_p
\]
Phương trình \eqref{eq:main} và $x \equiv 1 \pmod 3$, $a \nmid b$, $\vec{u}\perp\vec{v}$, $\widehat{ABC}=60^\circ$.

\subsection{Danh sách và bố cục}
\begin{tasks}(4)
\task $x=1$ \task $x=2$ \task $x=3$ \task $x=4$
\end{tasks}
\begin{description}
\item[Định lí Vi-ét] Tổng hai nghiệm $x_1+x_2=-\frac{b}{a}$.
\item[Hệ quả] Tích hai nghiệm.
\end{description}
\begin{enumerate}
\item Mức 1
  \begin{enumerate}
  \item Mức 2
    \begin{enumerate}
    \item Mức 3
    \end{enumerate}
  \end{enumerate}
\item Mức 1 tiếp
\end{enumerate}

\begin{tcolorbox}[colback=yellow!10, colframe=xanh, title=Ghi nhớ]
Công thức nghiệm: $x=\dfrac{-b\pm\sqrt{\Delta}}{2a}$.
\end{tcolorbox}

\begin{multicols}{2}
Cột trái với nhiều chữ để thấy hai cột được chia ra trong Word, có đủ nội dung để kiểm tra.
\columnbreak
Cột phải cũng có nội dung.
\end{multicols}

\begin{flushright}
\textit{Người ra đề}\\
\textbf{Nguyễn Văn A}
\end{flushright}

\begin{figure}[ht]
\centering
\includegraphics[width=0.5\textwidth]{hinh1.png}
\caption{Đồ thị hàm số}\label{fig:1}
\end{figure}
Hình~\ref{fig:1} cho thấy\ldots{} Trích dẫn \cite{sgk}.

\begin{tabularx}{\textwidth}{|X|c|}
\hline
\rowcolor{gray!20} Cột rộng & Hẹp \\ \hline
Nội dung dài dài dài & \cellcolor{green!20} 1 \\ \hline
\end{tabularx}

\begin{thebibliography}{9}
\bibitem{sgk} Bộ GD\&ĐT, \emph{Sách giáo khoa Toán 12}, NXB Giáo dục, 2024.
\end{thebibliography}

Ký tự đặc biệt: \% \$ \& \# \_ \{ \} \textbackslash{} \S\ \copyright{} 50\,\% --- dấu ``nháy'' và ~cách. Accent: \'e \`a \^o \"u \~n \c{c} \d{a} \h{a}.
Mã \verb|\frac{a}{b}| giữ nguyên.
\end{document}
